% ============================================================ broken symmetry and spin projection
\section{Broken symmetry}

\begin{frame}{Cheap methods for diradicals}
  \begin{paperclaim}
    ``relatively `cheap' methods such as unrestricted or broken symmetry formalisms suffer from a number of serious
    drawbacks, while rigorous methods \dots\ are generally not feasible for systems large enough to be practically
    valuable.'' (p.\ 443)
  \end{paperclaim}
  \begin{columns}[c]
    \column{0.55\textwidth}
    broken-symmetry (BS) determinant $\ket{a'\bar b'}$, \ $a'$ on centre $l$, $b'$ on $r$, overlap $s=\braket{a'}{b'}$
    {\large
    \begin{align*}
      \uncover<2->{\ev{\hat S^2}_{\text{UHF}}}&\uncover<2->{=S_z(S_z+1)+N_\beta-\textstyle\sum_{i\alpha,j\beta}|\braket ij|^2}\\
      \uncover<3->{\ev{\hat S^2}_{\text{BS}}}&\uncover<3->{=0+1-s^2}
    \end{align*}}
    \column{0.4\textwidth}
    \onslide<4->{\small $s=1$: closed shell, $\ev{\hat S^2}=0$\\ $s=0$: half singlet, half triplet, $\ev{\hat S^2}=1$}\\[4pt]
    \onslide<5->{\cgold{a single determinant cannot be the open-shell singlet} (Part III)}
  \end{columns}
\end{frame}

\begin{frame}{Yamaguchi's projection, derived}
  \begin{assume}
    two-centre Heisenberg model $\hat H=-2J\,\hat{\mathbf S}_1\!\cdot\!\hat{\mathbf S}_2$; BS state $=c_S\Psi_S+c_T\Psi_T$ ($M_S=0$)
  \end{assume}
  {\large
  \begin{align*}
    \stepno{1}\ \hat{\mathbf S}_1\!\cdot\!\hat{\mathbf S}_2&=\tfrac12\bigl[S(S+1)-\tfrac32\bigr]
      \ \Rightarrow\ E_T=-\tfrac J2,\ E_S=\tfrac{3J}2,\ E_S-E_T=2J\\
    \uncover<2->{\stepno{2}\ \ev{\hat S^2}_{\text{BS}}}&\uncover<2->{=2|c_T|^2,\qquad E_{\text{BS}}=|c_S|^2E_S+|c_T|^2E_T}\\
    \uncover<3->{\stepno{3}\ E_{\text{BS}}-E_T}&\uncover<3->{=|c_S|^2(E_S-E_T)=\bigl(1-\tfrac12\ev{\hat S^2}_{\text{BS}}\bigr)2J}\\
    \uncover<4->{\stepno{4}\ J}&\uncover<4->{=\dfrac{E_{\text{BS}}-E_{\text{HS}}}{\ev{\hat S^2}_{\text{HS}}-\ev{\hat S^2}_{\text{BS}}}\qquad(\ev{\hat S^2}_{\text{HS}}=2)}
  \end{align*}}
  \onslide<5->{\centering\small\cmuted{Lean: \texttt{yamaguchi}; sign conventions differ between sources, here $J<0$ is antiferromagnetic}\par}
\end{frame}

\begin{frame}{How good is it? The Hubbard dimer}{T14: self-consistent BS-UHF vs the exact singlet--triplet gap}
  \begin{columns}[c]
    \column{0.58\textwidth}
    \begin{tikzpicture}
      \begin{axis}[noirplot,width=7.6cm,height=5.4cm,xmin=2,xmax=16,ymin=-0.62,ymax=0,
          xlabel={$U/t$},ylabel={$J/t$},ymajorgrids,legend pos=south east]
        \addplot[noircyan,very thick] table[x=U,y=Jexact,col sep=comma]{data/yamaguchi.csv};
        \addlegendentry{exact $\tfrac14(U-\sqrt{U^2+16t^2})$}
        \only<2->{\addplot[noirgold,very thick,dashed] table[x=U,y=Jyam,col sep=comma]{data/yamaguchi.csv};}
        \only<2->{\addlegendentry{Yamaguchi from BS-UHF}}
      \end{axis}
    \end{tikzpicture}
    \column{0.38\textwidth}
    \begin{itemize}
      \item<2-> $U=3t$: error $-7.7\%$
      \item<3-> $U=12t$: $-0.1\%$
      \item<4-> good when the centres are weakly coupled ($\ev{\hat S^2}_{\text{BS}}\to1$), poor when covalent
    \end{itemize}
    \onslide<5->{\okpill\ \small ``serious drawbacks'' is fair}
  \end{columns}
\end{frame}
